% Generated draft: requires mathematical and source review.
\documentclass[11pt]{article}
\usepackage[margin=25mm]{geometry}
\usepackage{fontspec}
\setmainfont{Latin Modern Roman}
\usepackage{amsmath,amssymb,mathtools,unicode-math}
\setmathfont{Latin Modern Math}
\usepackage{enumitem,needspace,xurl,hyperref}
\setlength{\parindent}{0pt}
\setlength{\parskip}{5pt}
\setlength{\emergencystretch}{3em}
\newcommand{\N}{\mathbb{N}}
\newcommand{\Z}{\mathbb{Z}}
\newcommand{\Q}{\mathbb{Q}}
\newcommand{\R}{\mathbb{R}}
\newcommand{\C}{\mathbb{C}}
\newcommand{\F}{\mathbb{F}}
\newcommand{\A}{\mathbb{A}}
\newcommand{\PP}{\mathbb{P}}
\newcommand{\E}{\mathbb{E}}
\newcommand{\eps}{\varepsilon}
\newcommand{\dd}{\,\mathrm{d}}
\begin{document}
\section*{UnsolvedMath ID: 20001914}
\textbf{Source problem code:} catalogue entry
\subsection{Definitions and mathematical statement}
Let \(M\) be a smooth closed connected manifold and let \(\mathcal G\) be its space of smooth Riemannian metrics, equipped with a specified finite \(C^k\) topology, or with the \(C^\infty\) topology when explicitly requested. For \(g\in\mathcal G\), let \(S_gM\) be the unit tangent bundle and let \(\operatorname{Per}(g)\subset S_gM\) be the tangent vectors on periodic geodesic-flow orbits. Question 8.4.3 asks whether
\[
 \{g\in\mathcal G:\overline{\operatorname{Per}(g)}=S_gM\}
\]
is dense in the intended metric topology. The source asks for phase-space density, which is stronger than density of the images of closed geodesics in \(M\). Its following Section 8.5 heading introduces that different base-space subject without supplying a further question here.

For the aid's finite-resolution statement, assume \(\dim M\geq2\), fix a smooth reference metric \(g_0\) and a compatible distance \(d_0\) on the compact reference bundle \(S_{g_0}M\). Identify bundles by radial normalization
\[
 j_g:S_gM\to S_{g_0}M,\qquad j_g(v)=v/|v|_{g_0},
\]
and put \(P(g)=j_g(\operatorname{Per}(g))\). The reported result is
\[
 \forall m\in\N_{\geq1}\ \exists\mathcal O_m\subset\mathcal G:
 \quad\mathcal O_m\text{ is }C^1\text{-open and }C^1\text{-dense},
\]
\[
 g\in\mathcal O_m\quad\Longrightarrow\quad
 \forall z\in S_{g_0}M\ \exists v\in P(g):\ d_0(z,v)<1/m.
\]
Thus every prescribed finite resolution is attained on an open dense set, with the metric allowed to depend on that resolution. The argument is reported to combine a one-direction \(C^1\) closing theorem of Rifford, an orbit-preserving Ma\~n\'e perturbation of Rifford--Ruggiero to make the periodic orbit nondegenerate, fixed-point-index persistence under \(C^1\) metric perturbations, and a finite phase-space net. The precise hypotheses and references of these ingredients are not supplied.

This quantifier order does not by itself provide one smooth metric satisfying every resolution, or a dense set of such metrics. Smooth metrics with a fixed finite \(C^k\) topology are not a Baire space, so taking \(\bigcap_m\mathcal O_m\) has no automatic dense-intersection conclusion. Moreover \(C^1\)-density does not establish density in stronger \(C^k\) or \(C^\infty\) topologies.

The aid also reports a high-jet sufficient criterion. A precise Baire formulation is: take a countable basis \((U_j)\) of nonempty open subsets of \(S_{g_0}M\), and suppose each periodic hitting condition
\[
 \mathcal H_j=\{g:P(g)\cap U_j\ne\varnothing\}
\]
contains a \(C^\infty\)-open dense subset. Density here must allow perturbations inside every prescribed finite-jet neighborhood: for every smooth \(g\), finite derivative order and positive tolerance, there is a nearby smooth metric with the required persistent periodic hit. Since \(\mathcal G\) with its smooth Fréchet topology is Baire, these hypotheses yield a residual \(C^\infty\)-dense set with \(P(g)\) dense. This is a conditional logical criterion; the \(C^1\) perturbation results do not supply its missing high-jet control.

The original background, citing [Weinstein1970], reports examples with an open subset of the unit tangent bundle containing no periodic tangent vector, showing that exact density is not asserted for every metric. Its phrase ``any manifold'' is not dimension-qualified in the extraction and must be checked; in particular circle metrics have all unit-speed geodesics periodic. The finite-resolution theorem above is explicitly in dimension at least two and does not settle the original exact-density problem in the supplied account.
\subsection{Short English statement}
At any fixed phase-space resolution, the aid reports an open dense set in the \(C^1\) topology with sufficiently many periodic directions. This does not automatically give exact density for a single smooth metric, nor density in a stronger topology. A separate high-jet perturbation criterion would yield a smooth generic result.
\subsection{Sources}
\begin{itemize}
\item Permanent UnsolvedMath ID \(20001914\), exact public code catalogue entry; local registry \texttt{aim-geometry-notes.json}, \texttt{notes[251]}, canonical tag question.
\item AIM workshop ``Geodesics,'' \emph{Open Problems and Questions about Geodesics}, Question 8.4.3, extracted source item 8.4: \url{https://aimath.org/pastworkshops/geodesicsproblems.pdf}. The supplied verification note identifies the smooth-metric space and finite \(C^k\) topology. No PDF was accessed.
\item Original reference label [Weinstein1970]. The aid additionally names Rifford's one-direction \(C^1\) closing theorem and a Rifford--Ruggiero orbit-preserving Ma\~n\'e perturbation, without full bibliographic details. No original literature entry, proposer or year is supplied.
\item Research attempt 1, ``Finite-resolution density of periodic geodesic directions,'' labelled partial\_result, open at the run, proved and low novelty confidence; registry partially\_solved. The finite-resolution and high-jet-criterion descriptions are supplied research reports, not independent verification.
\end{itemize}
\end{document}
